\section{The exact finite expectation kernel}\label{sec:finite}

Throughout this section \(L=(m,w,W)\) is a finite law and
\(X,Y:\Fin(m)\to\Q\) are rational observables.  Every identity is understood
both as ordinary rational mathematics and as its denominator-cleared coded
form.

\begin{definition}[Expectation, moment, and event probability]\label{def:expect}
Define
\[
 \E_L[X]\defeq\frac{\sum_{i<m}w_iX_i}{W},\qquad
 \E_L[X^r]\defeq\frac{\sum_{i<m}w_iX_i^r}{W},
\]
and, for a decidable event \(A\),
\[
 \Pp_L(A)\defeq\E_L[\ind_A]
 =\frac{\sum_{i<m}w_i\ind_A(i)}{W}.
\]
\end{definition}

\begin{theorem}[Normalization]\label{thm:normalization}
\(\Pp_L(\Fin(m))=1\), \(\Pp_L(\varnothing)=0\), and
\(0\le\Pp_L(A)\le1\) for every event \(A\).
\end{theorem}

\begin{proof}
For the whole space the indicator is constantly one, so its weighted sum is
\(\sum_iw_i=W\).  For the empty event it is constantly zero.  Finally,
\(0\le w_i\ind_A(i)\le w_i\) pointwise; bounded-sum monotonicity gives
\(0\le\sum_iw_i\ind_A(i)\le W\), and division by the positive \(W\) preserves
order.  Internally one retains the cross-multiplied numerator inequalities.
\end{proof}

\begin{theorem}[Weighted-average formula]\label{thm:weighted-average}
If \(p_i=w_i/W\), then \(p_i\ge0\), \(\sum_i p_i=1\), and
\[
 \E_L[X]=\sum_{i<m}p_iX_i.
\]
Conversely every finite rational probability vector has such a natural-weight
presentation.
\end{theorem}

\begin{proof}
The forward statements are direct expansion.  For the converse, take the
product of all positive denominators as a common \(W\) and put \(w_i=Wp_i\).
The equation \(\sum p_i=1\) becomes \(\sum w_i=W\) after multiplication by
\(W\).
\end{proof}

\begin{theorem}[Convex-hull bounds]\label{thm:convex-hull}
Suppose rational numbers \(a,b\) satisfy
\(a\le X_i\le b\) for every positive-weight outcome.  Then
\[
 a\le\E[X]\le b.
\]
If \(\E[X]=a\), then \(X_i=a\) at every positive-weight outcome whenever
\(X_i\ge a\); similarly at the upper endpoint.
\end{theorem}

\begin{proof}
Pointwise, \(w_i a\le w_iX_i\le w_i b\).  Sum and use
\(\sum_iw_i=W\).  For the endpoint assertion write \(X_i-a\ge0\).  The equality
\(\E[X]=a\) says \(\sum_iw_i(X_i-a)=0\) after denominator clearing.  Every
nonnegative summand is therefore zero.  If \(w_i>0\), cancellation yields
\(X_i=a\).
\end{proof}

\begin{corollary}[Equiprobable arithmetic mean]\label{cor:equiprobable}
For equal positive weights,
\[
 \E[X]=\frac{X_0+\cdots+X_{m-1}}{m}.
\]
In particular, expectation generalizes the ordinary arithmetic mean.
\end{corollary}

\begin{proof}
Take \(w_i=1\) and \(W=m\) in \Cref{def:expect}.
\end{proof}

\subsection{Pushforwards and finite LOTUS}

\begin{definition}[Pushforward law]\label{def:pushforward}
Let \(f:\Fin(m)\to\Fin(n)\) be primitive recursive.  Its pushforward weights
are
\[
 v_j\defeq\sum_{\substack{i<m\\f(i)=j}}w_i.
\]
The total remains \(\sum_{j<n}v_j=W\).
\end{definition}

\begin{theorem}[Finite law of the unconscious statistician]\label{thm:finite-lotus}
For every rational \(g:\Fin(n)\to\Q\),
\[
 \E_L[g(f(i))]=\frac1W\sum_{j<n}v_jg(j).
\]
Equivalently, expectation after applying a function depends only on the
pushforward distribution.
\end{theorem}

\begin{proof}
Insert a decidable partition of unity:
\[
 \sum_{i<m}w_i g(f(i))
 =\sum_{i<m}\sum_{j<n}w_i\ind_{f(i)=j}g(j).
\]
For each \(i\), exactly one indicator is one, so the inner sum equals the
original term.  Apply finite Fubini and factor out \(g(j)\):
\[
 \sum_{j<n}g(j)\sum_{i<m}w_i\ind_{f(i)=j}
 =\sum_{j<n}v_jg(j).
\]
Every equality is a bounded-sum identity.  The claim that exactly one fiber
indicator is one is proved by decidable equality on \(\Fin(n)\).
\end{proof}

\begin{corollary}[Multidimensional LOTUS]\label{cor:multi-lotus}
For finite observables \(X_1,\ldots,X_r\) and a rational function \(g\) on
their finite joint range,
\[
 \E[g(X_1,\ldots,X_r)]
 =\sum_{x_1,\ldots,x_r}g(x_1,\ldots,x_r)
   \Pp(X_1=x_1,\ldots,X_r=x_r).
\]
\end{corollary}

\begin{proof}
Code the tuple by primitive-recursive pairing and apply
\Cref{thm:finite-lotus}.  Repeated finite Fubini converts the single paired
sum to the displayed nested sum.
\end{proof}

\subsection{Vector and matrix expectation}

\begin{theorem}[Componentwise expectation]\label{thm:componentwise}
Let \(X(i)=(X_1(i),\ldots,X_r(i))\in\Q^r\).  Define expectation of a vector by
its finite weighted sum.  Then
\[
 \E[X]=(\E[X_1],\ldots,\E[X_r]).
\]
For a rational matrix-valued observable, expectation is entrywise.  Rational
matrix multiplication by a fixed matrix commutes with expectation.
\end{theorem}

\begin{proof}
A coded vector or matrix is a finite tuple of rational codes.  Projection onto
a coordinate commutes definitionally with the bounded weighted sum.  For a
fixed matrix \(A\), each coordinate of \(AX\) is a finite linear combination
of coordinates of \(X\); distributivity of bounded sums proves
\(\E[AX]=A\E[X]\).  No higher-type object is needed: the dimension and all
coordinate access functions are natural-number parameters.
\end{proof}

\subsection{Finite distributions as normal forms}

\begin{proposition}[Histogram normal form]\label{prop:histogram}
Every finite rational observable has a primitive-recursively computed finite
histogram \((x_j,v_j)_{j<n}\) such that the \(x_j\) are pairwise distinct,
\(v_j\) is the total weight of outcomes with value \(x_j\), and
\[
 \E[X]=\frac1W\sum_{j<n}v_jx_j.
\]
The histogram may be obtained without sorting by retaining one slot per input
outcome and allowing zero duplicate slots.
\end{proposition}

\begin{proof}
Use the pushforward map from outcome indices to the least earlier index with
the same rational value; bounded search computes this representative.
Alternatively, define
\(v_j=w_j\ind_{j\text{ is first}}+\sum_{i>j}w_i\ind_{X_i=X_j}\ind_{j\text{ is first}}\).
Both are primitive-recursive.  Fiber regrouping gives the expectation identity.
The zero-padding option avoids any need for a variable-length quotient type.
\end{proof}

\begin{remark}
The histogram form is useful when translating textbook notation
\(\sum_x x\Pp(X=x)\).  The sample-space form is usually simpler for formal
proofs, because no proof of distinctness or enumeration completeness is
required.
\end{remark}
